\section{Results and Discussion}
\label{sec:results}

All results in this section come from the single specification set out in
Section~\ref{sec:method}. No parameter was selected by reference to the evaluation
period, and every specification examined is reported, including those that are
unfavourable to the framework.

\subsection{Overview of Experimental Results}

The implemented decision-theoretic framework provides evidence for the efficacy of
regime-aware portfolio construction, and evidence against two of the mechanisms it was
built around. Across the 4{,}689 trading sessions from 7~January 2008 to 26~August 2026,
encompassing diverse market conditions, the regime-based strategy achieved a Sharpe ratio
of 0.878 compared with the equal-weight benchmark's 0.628, a 40~per~cent improvement in
risk-adjusted return. More significantly, maximum drawdown was reduced from $-36.08$ to
$-18.98$ per cent, highlighting the defensive capability of the regime-adaptive approach
during market stress. Table~\ref{tab:main} reports the primary comparison.

\begin{table}[!ht]
\centering
\caption{Primary results, ten-asset universe, January 2008 to August 2026. Sharpe
ratios use a zero risk-free rate and are computed as compound annual growth divided by
annualised volatility. Benchmarks are rebalanced daily and charged no transaction costs.}
\label{tab:main}
\footnotesize
\begin{tabular}{p{4.4cm}rrrrr}
\hline
& \textbf{CAGR} & \textbf{Volatility} & \textbf{Sharpe} & \textbf{Max DD} & \textbf{Calmar} \\
\hline
Regime-conditional strategy & 5.65\% & 6.44\% & 0.878 & $-18.98$\% & 0.298 \\
Equal-weight benchmark & 8.45\% & 13.45\% & 0.628 & $-36.08$\% & 0.234 \\
Inverse-volatility benchmark & 7.46\% & 8.78\% & 0.849 & $-24.09$\% & 0.309 \\
Equal-weight at strategy volatility & 4.19\% & 6.44\% & 0.651 & $-18.69$\% & 0.224 \\
\hline
\end{tabular}
\end{table}

Three readings of Table~\ref{tab:main} are warranted, and they do not all favour the
framework.

Against the equal-weighted benchmark the framework reduces maximum drawdown by 47 per
cent while earning a lower compound return, 5.65 per cent against 8.45 per cent per
annum. It is not a return-enhancing strategy and is not presented as one; it is a
risk-reduction strategy whose return sacrifice is smaller than its risk reduction.

Against the volatility-matched equal-weighted benchmark, which removes the mechanical
advantage a lower-risk portfolio enjoys in any ratio comparison, the framework retains a
Sharpe advantage of 0.878 against 0.651 but loses on maximum drawdown, $-18.98$ per cent
against $-18.69$ per cent. Deleveraging a diversified benchmark is therefore a competitive
way to reduce worst-case loss, and the framework's drawdown advantage over the unlevered
benchmark should not be read as evidence that it dominates every risk-reduced
alternative.

Against the inverse-volatility benchmark the comparison is mixed: a higher Sharpe ratio,
0.878 against 0.849, and a materially smaller maximum drawdown, $-18.98$ per cent against
$-24.09$ per cent, but a lower Calmar ratio, 0.298 against 0.309, because
inverse-volatility weighting compounds faster. Elementary risk weighting remains a strong
comparator on a diversified multi-asset universe, and the margin over it is narrow enough
to read as a tie on risk-adjusted return rather than a win.

\subsection{Detailed Analysis of Visualizations and Performance Metrics}

\subsubsection{Cumulative Returns Comparison}

\begin{figure*}[t]
\centering
\includegraphics[width=\textwidth]{images/fig_cumret.png}
\caption{Cumulative Returns Comparison: cumulative growth of one unit of capital,
logarithmic scale}
\label{fig:cumret}
\end{figure*}

The comparative analysis of cumulative returns, presented in Figure~\ref{fig:cumret},
delineates a performance profile characterised by distinct phases of divergence and
convergence relative to the passive benchmark. The regime-aware portfolio exhibits
several defining characteristics over the observed period:

\begin{itemize}
\item The strategy separates upward during periods of significant market dislocation,
notably 2008, 2011, the volatility surge in early 2020 and the inflationary selloff of
2022.
\item A critical observation is the strategy's effective mitigation of losses during
sustained market downturns, indicating a significantly reduced downside capture: 0.38,
against an upside capture of 0.41.
\item During extended bullish phases---most conspicuously 2013, 2017, 2021 and
2023---the equity curve tracks below the benchmark, so the strategy delivers lower
absolute returns while exhibiting a lower volatility signature, 6.44 against 13.45 per
cent.
\end{itemize}

\textbf{Strategic Implications:}

The observed trajectory provides evidence for the regime detection module's operational
efficacy. The system's identification of the transition into a high-volatility state in
early 2020 triggered a timely de-risking of the portfolio, constraining the drawdown of
that episode to $-6.29$ per cent against $-21.28$ per cent for the benchmark, and the
subsequent gradual ascent in the portfolio's value line indicates a methodical
re-allocation of risk exposure as the regime normalised, avoiding premature or delayed
re-entry. The cost is equally visible and is not concealed: the framework forgoes 2.80
percentage points of annual compound return over the sample.

\textbf{Theoretical Alignment:}

These results substantiate one half of the foundational thesis of this research and
refute the other. The performance pattern confirms that the total risk a portfolio should
carry is conditional upon the underlying market state, and the capture ratio of 0.93
locates the source precisely: the framework participates less in both directions in
roughly equal proportion, so its advantage comes from carrying less risk in adverse
states rather than from directional skill. The dynamic fusion mechanism did not operate as
theorised---Section~\ref{sec:ic} shows the conditional informativeness its weighting
schedule presumes is not present in the data, and Section~\ref{sec:ablation} shows that
removing the fusion layer does not harm the portfolio.

\subsubsection{Portfolio Drawdown}

\begin{figure*}[t]
\centering
\includegraphics[width=\textwidth]{images/fig_drawdown.png}
\caption{Portfolio Drawdown: drawdown from running peak}
\label{fig:drawdown}
\end{figure*}

The comparative drawdown analysis, illustrated in Figure~\ref{fig:drawdown}, provides
evidence of the regime-based portfolio's capital preservation attributes. The visual data
reveal a consistent pattern of defensive characteristics:

\begin{itemize}
\item The strategy registers a significantly shallower maximum drawdown of $-18.98$ per
cent, against the benchmark's $-36.08$ per cent nadir; the benchmark reaches $-36$ per
cent in 2008 and $-25$ per cent in 2022, while the framework's deepest excursions are
approximately half as large.
\item Following drawdown troughs, the regime-based portfolio demonstrates a swifter
recovery trajectory, spending less time underwater, since a shallower loss requires a
smaller subsequent gain to restore the previous peak.
\item The portfolio experiences briefer episodes of significant capital impairment,
indicating more timely risk mitigation. On the worst 5~per~cent of benchmark sessions
($n = 234$) it returns $-0.59$ per cent on average against $-2.01$ per cent for the
benchmark, a cushion of 1.42 percentage points per session.
\end{itemize}

Table~\ref{tab:episodes} reports cumulative returns over four stress episodes chosen
before the analysis on the basis of documented market history rather than by inspection
of the results.

\begin{table}[!ht]
\centering
\caption{Cumulative return over stress episodes. Differences are stated in percentage
points (pp).}
\label{tab:episodes}
\footnotesize
\begin{tabular}{p{3.6cm}rrrr}
\hline
\textbf{Episode} & \textbf{Sessions} & \textbf{Strategy} & \textbf{Equal-weight} & \textbf{Difference} \\
\hline
Global financial crisis, Sep 2008--Mar 2009 & 130 & $-7.05$\% & $-29.78$\% & $+22.74$ pp \\
Euro area crisis, Jul--Oct 2011 & 65 & $+3.41$\% & $-7.24$\% & $+10.65$ pp \\
COVID-19 crash, Feb--Mar 2020 & 24 & $-6.29$\% & $-21.28$\% & $+14.99$ pp \\
Inflation bear market, Jan--Oct 2022 & 196 & $-16.76$\% & $-24.49$\% & $+7.73$ pp \\
\hline
\end{tabular}
\end{table}

\textbf{Risk Management Effectiveness:}

The observed drawdown profile validates the regime detection system as an early-warning
mechanism for deteriorating market states. By signalling a transition towards a
high-volatility state, the framework enabled proactive de-risking, and it outperforms in
all four episodes by between 7.7 and 22.7 percentage points, the widest margin being the
global financial crisis, the longest and deepest of the four. The margin is narrowest in
the 2022 bear market, consistent with the design: a gradual, correlated decline across
both equities and bonds offers less scope for a de-risking rule than an abrupt volatility
shock does. All four episodes predate the holdout boundary; the deepest benchmark decline
after it, the 34-session decline from 20~February to 8~April~2025, produced an advantage
of 5.42 percentage points.

\textbf{Theoretical Contribution:}

This outcome substantiates the value of conditioning total risk on a state variable, and
bounds that claim in one respect a comparison against an unlevered benchmark would
conceal. Against the volatility-matched benchmark the drawdown advantage disappears,
$-18.98$ against $-18.69$ per cent, and Section~\ref{sec:robust} shows this holds in every
one of the eighteen specifications examined. The framework therefore reduces drawdown
relative to an unlevered diversified benchmark, not relative to every portfolio held at
the same volatility; reporting only the first comparison, as the conventional presentation
in this literature does, would overstate the result.

\subsubsection{Market Regime Detection}
\label{sec:regimefit}

\begin{figure*}[t]
\centering
\includegraphics[width=0.92\textwidth]{images/fig_regimes.png}
\caption{Market Regime Detection: identified market states, with the underlying
volatility and momentum features}
\label{fig:regimes}
\end{figure*}

The temporal sequence of identified market states, presented in
Figure~\ref{fig:regimes}, delineates three economically distinct market environments
characterised by their unique volatility and momentum signatures:

\begin{itemize}
\item \textbf{Calm.} This state is characterised by favourable market conditions,
combining low volatility with sustained positive momentum, typically associated with
bullish trending markets: 12.22 per cent annualised volatility and $+16.04$ per cent
annualised momentum.
\item \textbf{Transitional.} Representing a neutral state, this regime captures market
conditions that lack a definitive directional trend or persistent volatility
characteristics: 20.04 per cent volatility and $-1.25$ per cent momentum.
\item \textbf{Stress.} This state corresponds to periods of elevated market stress,
marked by high volatility and negative momentum, at 47.07 per cent and $-15.97$ per cent
respectively. Its incidence is clustered around known macroeconomic disruptions,
specifically the 2008 global financial crisis, the 2011 sovereign debt episode, the
initial COVID-19 market crisis and the 2022 inflationary selloff.
\end{itemize}

Table~\ref{tab:regimefit} reports the model's order selection statistics and the
properties of the three-state solution.

\begin{table}[!ht]
\centering
\caption{Regime model diagnostics. Order selection is computed on the full feature
history; state characteristics and the transition matrix come from the three-component
solution. Out-of-sample shares are the states actually assigned by the expanding-window
refits used in the backtest.}
\label{tab:regimefit}
\footnotesize
\begin{tabular}{p{3.2cm}rrrrr}
\hline
\multicolumn{6}{l}{\textit{Order selection}} \\
\hline
\textbf{Components} & 2 & 3 & 4 & 5 & 6 \\
AIC & 21944 & 20785 & 20571 & 20478 & 20372 \\
BIC & 22016 & 20896 & 20720 & 20667 & 20599 \\
\hline
\multicolumn{6}{l}{\textit{Three-state solution}} \\
\hline
& \textbf{Calm} & \textbf{Transitional} & \textbf{Stress} & & \\
Annualised volatility & 12.22\% & 20.04\% & 47.07\% & & \\
Annualised momentum & $+16.04$\% & $-1.25$\% & $-15.97$\% & & \\
Share of sessions, in sample & 63.6\% & 31.2\% & 5.2\% & & \\
Share of sessions, out of sample & 71.4\% & 25.5\% & 3.1\% & & \\
Mean run length (sessions) & 25.3 & 11.1 & 17.1 & & \\
Self-transition probability & 0.961 & 0.910 & 0.941 & & \\
\hline
\multicolumn{6}{l}{\textit{Realised performance within each out-of-sample state}} \\
\hline
& \textbf{Calm} & \textbf{Transitional} & \textbf{Stress} & & \\
Sessions & 3346 & 1196 & 147 & & \\
Strategy Sharpe & 1.144 & 0.196 & 1.039 & & \\
Equal-weight Sharpe & 0.916 & 0.561 & 0.356 & & \\
\hline
\end{tabular}
\end{table}

\textbf{Temporal Regime Characteristics:}

An examination of regime persistence reveals distinct behavioural patterns. All three
states persist, with self-transition probabilities between 0.910 and 0.961. The
transitional state exhibits the briefest durations, a mean run length of 11.1 sessions,
indicating its nature as a short-lived bridge between the other two. In contrast, the
calm state demonstrates the greatest persistence, enduring 25.3 sessions on average, with
the stress state intermediate at 17.1. Transitions between these states are constrained:
the transition matrix is close to tridiagonal, and the estimated probability of moving
directly between the calm and the stress state, in either direction, is zero. The
ordering recovered by sorting components on volatility is therefore economically
meaningful rather than arbitrary.

\textbf{Methodological Validation:}

The output from the Gaussian mixture model supports the regime detection methodology.
The model's identification of stress periods aligns with historically documented stress
events in 2008, 2011, 2020 and 2022. This congruence between statistically identified
states and ex-post economic reality indicates the framework's capacity to discriminate
between structurally different market environments, a foundational requirement for a
regime-dependent investment strategy. The states also separate sharply on a variable
withheld from their construction, as Section~\ref{sec:sentval} reports.

\textbf{Order Selection:}

The order selection statistics are reported rather than used. Both AIC and BIC continue
to fall over the range examined, from 21944 and 22016 at two components to 20372 and
20599 at six, so neither criterion selects an interior optimum and neither supports the
choice of three components. This behaviour is expected: a mixture fitted to persistent,
heavy-tailed financial features will absorb tail observations into additional components
indefinitely, and the criteria reward it for doing so. The choice of three is made instead
on persistence and interpretability, and is reported as a modelling judgement rather than
as a selected optimum.

\textbf{Performance Within Each State:}

The final panel of Table~\ref{tab:regimefit} shows where the framework's advantage
arises. In the stress state the framework attains a Sharpe ratio of 1.039 against 0.356
for the benchmark, and in the calm state 1.144 against 0.916. It underperforms materially
only in the transitional state, 0.196 against 0.561, which is consistent with the 2021
result in Table~\ref{tab:years}: a state that is genuinely ambiguous causes the risk
budget to reduce exposure in advances that turn out to continue. The mechanism is
therefore effective at the two ends of the volatility range and costly in the middle,
which is an intelligible property of a three-state partition rather than a defect in it.

\subsubsection{Portfolio Allocation Over Time}
\label{sec:alloc}
\label{sec:cost}

\begin{figure*}[t]
\centering
\includegraphics[width=0.92\textwidth]{images/fig_alloc.png}
\caption{Portfolio Allocation: portfolio weights over time, including the cash residual}
\label{fig:alloc}
\end{figure*}

\begin{figure*}[t]
\centering
\includegraphics[width=\textwidth]{images/fig_exposure.png}
\caption{Invested exposure over time, with transitional and stress regimes shaded}
\label{fig:exposure}
\end{figure*}

The temporal evolution of portfolio weights, illustrated in Figure~\ref{fig:alloc},
provides a clear representation of the system's adaptive response to changing market
states, and Figure~\ref{fig:exposure} displays the total exposure it carries. The
allocation dynamics reveal a systematic rebalancing process:

\begin{itemize}
\item During periods identified as the stress state, capital was systematically withdrawn
from the risky sleeve toward cash, with invested exposure falling to a minimum of 50 per
cent.
\item Under transitional or calm market regimes, the portfolio maintained a balanced
exposure across multiple asset classes, avoiding significant concentration risk: the
effective number of assets never falls below 4.60.
\item The composition of the risky sleeve itself is set by the optimiser subject to the
25 per cent position cap, which binds on average, so the framework expresses its market
view through total exposure rather than through concentration.
\end{itemize}

This entire process is characterised by a rules-based systematicity, entirely devoid of
discretionary influence: the exposure at every rebalance is a deterministic function of
the regime label, and the composition a deterministic function of the fused score and the
covariance estimate.

Table~\ref{tab:turnover} summarises these characteristics together with the sensitivity
of performance to the assumed transaction cost.

\begin{table}[!ht]
\centering
\caption{Allocation characteristics and sensitivity to the assumed transaction cost.
Effective number of assets is the reciprocal of the Herfindahl index of sleeve weights.
Annual turnover is one-way, as a multiple of portfolio value.}
\label{tab:turnover}
\footnotesize
\begin{tabular}{p{4.8cm}rrrr}
\hline
\multicolumn{5}{l}{\textit{Allocation characteristics, primary specification}} \\
\hline
Mean largest weight & \multicolumn{4}{r}{0.250} \\
Mean effective number of assets & \multicolumn{4}{r}{5.40} \\
Minimum effective number of assets & \multicolumn{4}{r}{4.60} \\
Mean invested exposure & \multicolumn{4}{r}{92.1\%} \\
Minimum invested exposure & \multicolumn{4}{r}{50.0\%} \\
Annual one-way turnover & \multicolumn{4}{r}{0.96$\times$} \\
\hline
\multicolumn{5}{l}{\textit{Sensitivity to assumed cost}} \\
\hline
\textbf{Assumed cost} & \textbf{CAGR} & \textbf{Sharpe} & \textbf{Max DD} & \textbf{Turnover} \\
\hline
0 basis points & 4.90\% & 0.773 & $-19.83$\% & 4.04$\times$ \\
5 basis points & 5.66\% & 0.893 & $-19.54$\% & 1.14$\times$ \\
10 basis points (primary) & 5.65\% & 0.878 & $-18.98$\% & 0.96$\times$ \\
25 basis points & 5.47\% & 0.846 & $-19.30$\% & 0.93$\times$ \\
\hline
\end{tabular}
\end{table}

\textbf{Decision-Theoretic Implementation:}

The allocation is diversified and stable: annual one-way turnover of 0.96 times portfolio
value is equivalent to approximately 8.0 per cent of the portfolio traded at each monthly
rebalance. The observed smooth transitions between asset weights, rather than abrupt,
discontinuous shifts, reflect the integral role of the transaction cost-aware optimisation
function, which balances responding to new predictive signals against maintaining
implementation efficiency. The absence of extreme allocation swings throughout the
backtested period is evidence for the practical feasibility of the proposed framework in a
live implementation context.

\textbf{Cost Sensitivity:}

Performance is not monotone in the assumed cost. At zero assumed cost the Sharpe ratio is
0.773, below the 0.878 obtained at 10 basis points, even though no costs are deducted.
The reason is that the assumed cost enters the objective as a turnover penalty as well as
entering the return as a deduction. With the penalty removed, annual turnover rises from
0.96 to 4.04 times portfolio value as the optimiser chases small differences in the
estimated covariance matrix, and the resulting weight instability costs more than the
saved commissions are worth. The turnover penalty is therefore functioning as a
regulariser on the weight vector, and the framework is robust across the plausible cost
range of 5 to 25 basis points, over which the Sharpe ratio varies only between 0.846 and
0.893.

\textbf{A Correction to the Submitted Specification:}

This subsection also resolves a defect in the framework as originally specified. In that
version the fused score entered the objective without the information-ratio scaling of
Equation~(\ref{eq:alpha}), so the linear term was expressed in the raw units of the
underlying indicators. On-balance volume takes values of order $10^{8}$ to $10^{9}$,
while annualised variances are of order $10^{-2}$. The linear term consequently dominated
the quadratic risk term by roughly ten orders of magnitude, the risk model became
irrelevant to the solution, and the optimiser returned a corner solution holding a single
asset on every rebalance date---an effective number of assets of exactly 1.00, with one
holding persisting for 1{,}364 of 1{,}421 sessions. The scaling in
Equation~(\ref{eq:alpha}) removes the defect by placing the linear term on a return
scale, and Table~\ref{tab:turnover} reports the corrected behaviour. The failure was one
of units, not of the optimiser or of the risk model.

\subsubsection{Signal Strength Dynamics}
\label{sec:ic}

\begin{figure*}[t]
\centering
\includegraphics[width=\textwidth]{images/fig_signal.png}
\caption{Signal Strength Dynamics: mean absolute cross-sectional score of each composite,
21-session moving average}
\label{fig:signal}
\end{figure*}

The temporal evolution of composite signal strength, depicted in
Figure~\ref{fig:signal}, reveals a distinct pattern of variation. The intensity of the
aggregated predictive output is not static but fluctuates in a manner closely linked to
market phases:

\begin{itemize}
\item The most pronounced signal magnitudes are observed proximate to major market
turning points, which reflects the model's sensitivity to shifts in the underlying
volatility and momentum features.
\item During periods characterised by well-established directional price trends, signal
strength maintains a moderate, consistent level.
\item In calm, range-bound market conditions the composite registers its lowest
amplitude, reflecting a relative absence of cross-sectional dispersion.
\end{itemize}

Variation in amplitude is a property of the signal, however, and not evidence that the
signal is informative. Table~\ref{tab:ic} measures the predictive content of the two
composites directly, as the cross-sectional Spearman information coefficient between the
score on session $t$ and the return on session $t+1$.

\begin{table}[!ht]
\centering
\caption{Information coefficients against next-session cross-sectional returns.
Newey--West $t$-statistics with ten lags in parentheses.}
\label{tab:ic}
\footnotesize
\begin{tabular}{p{3.8cm}rrr}
\hline
& \textbf{Technical} & \textbf{Sentiment} & \textbf{Fused} \\
\hline
All sessions ($n=4{,}669$) & $-0.030$ ($-2.14$) & $-0.026$ ($-2.26$) & $-0.033$ ($-2.42$) \\
Calm regime ($n=3{,}339$) & $-0.042$ ($-2.48$) & $-0.040$ ($-3.02$) & --- \\
Transitional regime ($n=1{,}183$) & $-0.012$ ($-0.48$) & $+0.007$ ($+0.31$) & --- \\
Stress regime ($n=147$) & $+0.087$ ($+1.30$) & $+0.040$ ($+0.66$) & --- \\
\hline
2008--2023 ($n=4{,}025$) & $-0.035$ ($-2.31$) & $-0.032$ ($-2.67$) & $-0.039$ ($-2.73$) \\
2024--2026 holdout ($n=644$) & $+0.001$ ($+0.03$) & $+0.013$ ($+0.37$) & $+0.009$ ($+0.23$) \\
\hline
\end{tabular}
\end{table}

Over the full sample the information coefficients are negative and statistically
significant for both composites and for their fusion. A composite whose cross-sectional
score is negatively related to next-session returns cannot improve an allocation that
uses it in the stated direction, and the ablation in Table~\ref{tab:ablation} is the
portfolio-level consequence of the numbers in Table~\ref{tab:ic}.

\textbf{Signal-Regime Relationship:}

A granular analysis reveals a differentiated behaviour between signal types across
environments, though not the difference the framework was built to exploit. The regime
breakdown does not support the fusion hypothesis. The weighting schedule
assigns greater weight to the sentiment composite in the stress regime on the hypothesis
that it becomes relatively more informative there. The point estimates run the other way:
in the stress regime the technical composite reads higher than the sentiment composite,
and in the calm regime the two are indistinguishable. The contrast between calm and
stress is not statistically significant ($t = -1.46$, $p = 0.145$), so the correct
conclusion is that the data provide no evidence for state-dependent relative
informativeness in the direction the schedule assumes, rather than evidence for the
reverse. Either way the weighting schedule has no measured basis.

\textbf{Why the Negative Sign Is Not an Opportunity:}

The last two rows of Table~\ref{tab:ic} bear on how the negative coefficients should be
read. Splitting at the holdout boundary, the fused coefficient is $-0.039$ ($t = -2.73$)
over 2008--2023 and $+0.009$ ($t = 0.23$) over 2024--2026. The negative sign does not
persist. This matters because a stable negative information coefficient would be an
exploitable signal in inverted form, and the framework could be salvaged by reversing it.
The holdout rules that out: the in-sample negativity is a property of that era, not a
durable cross-sectional reversal effect, and inverting the composites would have bought
nothing after 2023. The honest reading is that these composites carry no reliable
directional information in either direction.

Two qualifications are owed. The stress regime contains only 147 sessions, so its
coefficients are imprecisely estimated and the positive technical reading should not be
over-interpreted. And the full-sample negative coefficients are consistent with
short-horizon cross-sectional reversal in a universe of broad index funds, which is a
known effect and not evidence of an error in construction, though the holdout shows it is
not stable enough to trade. The conclusion is bounded accordingly: these composites, as
specified, do not carry usable directional information in this universe at this horizon.

\subsubsection{Rolling Sharpe Ratio}

\begin{figure*}[t]
\centering
\includegraphics[width=\textwidth]{images/fig_rollsharpe.png}
\caption{Rolling Sharpe Ratio: rolling 252-session Sharpe ratio}
\label{fig:rollsharpe}
\end{figure*}

The 252-session rolling Sharpe ratio, presented in Figure~\ref{fig:rollsharpe}, shows
the regime-aware strategy's risk-adjusted performance across varying market conditions.
The trajectory is characterised by several defining features:

\begin{itemize}
\item The ratio maintains a positive value throughout the majority of the observed
period, in 81.9 per cent of windows, against 85.8 per cent for the benchmark.
\item Following periods of compression the metric reverts to its positive trend, and does
so from a shallower trough, since the framework's drawdowns are approximately half the
depth of the benchmark's.
\item It exceeds the benchmark in 53.4 per cent of windows, so the advantage is present
in slightly more windows than not rather than sustained throughout.
\end{itemize}

Table~\ref{tab:years} reports every calendar year in the sample, with both wins and
losses shown.

\begin{table}[!ht]
\centering
\caption{Calendar year performance. Sharpe and maximum drawdown computed within each
year. 2026 covers January to August only; the sample ends on 26~August 2026.}
\label{tab:years}
\footnotesize
\begin{tabular}{rrrrrcc}
\hline
\textbf{Year} & \textbf{Strategy} & \textbf{Equal-weight} & \textbf{Strat. Sharpe} & \textbf{EW Sharpe} & \textbf{Sharpe win} & \textbf{DD win} \\
\hline
2008 & $-2.28$\% & $-16.32$\% & $-0.33$ & $-0.47$ & yes & yes \\
2009 & $+4.36$\% & $+24.91$\% & $0.60$ & $1.18$ & no & yes \\
2010 & $+13.75$\% & $+18.31$\% & $2.05$ & $1.40$ & yes & yes \\
2011 & $+10.36$\% & $+5.75$\% & $1.73$ & $0.45$ & yes & yes \\
2012 & $+8.91$\% & $+13.49$\% & $1.96$ & $1.58$ & yes & yes \\
2013 & $-0.58$\% & $+6.17$\% & $-0.06$ & $0.72$ & no & yes \\
2014 & $+13.03$\% & $+9.97$\% & $2.94$ & $1.49$ & yes & yes \\
2015 & $-0.48$\% & $-1.56$\% & $-0.05$ & $-0.13$ & yes & yes \\
2016 & $+4.85$\% & $+8.49$\% & $0.94$ & $0.97$ & no & yes \\
2017 & $+10.92$\% & $+16.59$\% & $2.65$ & $3.15$ & no & yes \\
2018 & $-3.51$\% & $-5.26$\% & $-0.78$ & $-0.52$ & no & yes \\
2019 & $+15.72$\% & $+22.55$\% & $2.84$ & $3.02$ & no & yes \\
2020 & $+9.63$\% & $+19.32$\% & $1.23$ & $0.96$ & yes & yes \\
2021 & $-0.78$\% & $+10.15$\% & $-0.08$ & $1.12$ & no & no \\
2022 & $-12.99$\% & $-19.30$\% & $-1.55$ & $-1.23$ & no & yes \\
2023 & $+6.83$\% & $+16.53$\% & $0.90$ & $1.52$ & no & yes \\
2024 & $+9.38$\% & $+9.58$\% & $1.25$ & $1.00$ & yes & yes \\
2025 & $+18.53$\% & $+20.17$\% & $2.24$ & $1.67$ & yes & yes \\
2026 & $+4.80$\% & $+10.93$\% & $0.96$ & $1.34$ & no & yes \\
\hline
\end{tabular}
\end{table}

The framework wins on Sharpe ratio in 9 of 19 years and on drawdown in 18 of 19. This
distribution is the honest summary of the strategy: it loses on risk-adjusted return in
most individual years, particularly strong advancing years such as 2009, 2017, 2021 and
2023, and recovers the deficit in the minority of years when the benchmark suffers. The
three years after the holdout boundary follow the same pattern: drawdown wins in all
three, Sharpe wins in two.

Two years merit comment. In 2021 the framework loses on both measures, the only such
year, because the mixture model spent much of that year in the transitional state during
what turned out to be a sustained advance. In 2018 the framework loses on Sharpe despite
a smaller loss in absolute terms, an arithmetic consequence of a smaller negative return
divided by a much smaller volatility.

\textbf{Episodic Rather Than Persistent:}

The rolling window analysis bears on the question of performance persistence, and answers
it against the framework. The strategy's risk-adjusted return profile does not
consistently surpass that of a passive benchmark; it surpasses it in slightly more windows
than not, and the excess returns are clustered in the specific conditions catalogued in
Table~\ref{tab:episodes}. The advantage is therefore episodic rather than persistent,
which is what Table~\ref{tab:years} leads one to expect and what the stress-episode
results describe from the other direction. A reader evaluating this
framework should expect long stretches of underperformance punctuated by periods in which
the accumulated deficit is recovered, and should size the position accordingly.

\subsection{Robustness and Attribution}

The subsections that follow separate the framework into its components, test it outside
the sample on which it was specified, and report the inference that bounds every claim
made above.

\subsubsection{Out-of-Sample Holdout, 2024--2026}
\label{sec:holdout}

Every specification decision in this study---the universe, the feature set, the number of
regimes, the risk budget, the rebalancing interval and the cost assumption---was fixed
before any session after December~2023 was examined. The 665 sessions from January~2024
to August~2026 therefore function as a holdout rather than as a further robustness check
on the same data. Table~\ref{tab:holdout} reports both blocks side by side.

\begin{table}[!ht]
\centering
\caption{Performance before and after the holdout boundary. The two blocks are disjoint;
the full-sample results of Table~\ref{tab:main} pool them.}
\label{tab:holdout}
\footnotesize
\begin{tabular}{p{4.9cm}rrrr}
\hline
& \multicolumn{2}{c}{\textbf{2008--2023} ($n=4{,}024$)} & \multicolumn{2}{c}{\textbf{2024--2026} ($n=665$)} \\
& \textbf{Sharpe} & \textbf{Max DD} & \textbf{Sharpe} & \textbf{Max DD} \\
\hline
Regime-conditional strategy & 0.738 & $-18.98$\% & 1.614 & $-5.48$\% \\
Equal-weight benchmark & 0.532 & $-36.08$\% & 1.391 & $-10.13$\% \\
Inverse-volatility benchmark & 0.764 & $-24.09$\% & 1.397 & $-6.70$\% \\
Equal-weight at strategy volatility & 0.560 & $-17.66$\% & 1.378 & $-7.04$\% \\
No regime timing, full exposure & 0.785 & $-22.53$\% & 1.710 & $-6.47$\% \\
\hline
\end{tabular}
\end{table}

In the holdout the framework records a higher Sharpe ratio than all three benchmarks,
1.614 against 1.391, 1.397 and 1.378, and a smaller maximum drawdown than all three,
$-5.48$ per cent against $-10.13$, $-6.70$ and $-7.04$ per cent. Two of those comparisons
reverse the in-sample ordering: the inverse-volatility benchmark and the
volatility-matched equal-weighted benchmark both beat the framework on Sharpe ratio
before 2024 and both trail it afterwards. The deepest equal-weighted drawdown in the
holdout, the 34-session decline from 20~February to 8~April~2025, returned $-10.05$\% for
the benchmark against $-4.63$\% for the framework, an advantage of 5.42 percentage points
that is consistent with the four earlier episodes in Table~\ref{tab:episodes}.

Three qualifications limit what this establishes. First, 665 sessions are too few to
settle anything on their own: the bootstrap interval for the Sharpe difference in the
holdout is $[-0.352, 0.742]$, wider than the full-sample interval and comfortably
containing zero, and the framework again earns less in absolute terms, giving up 3.07
percentage points of annual return ($p = 0.348$). Second, the holdout contains no episode
comparable to 2008 or 2020; every series in it posts a Sharpe ratio above 1.3, so the
period is benign and the relevant information is the ordering rather than the level.
Third, the ablation ordering is unchanged---the variant with no regime timing at full
exposure again records the highest Sharpe ratio, 1.710---so the holdout does not rescue
any claim that Section~\ref{sec:ablation} rejects.

What the holdout does establish is narrower and still worth stating. The framework's
defensive behaviour is not an artefact of the sample on which it was specified. It was
built without reference to these sessions, and on them it continues to reduce drawdown
against every benchmark considered while remaining competitive on risk-adjusted return.

\subsubsection{Robustness Across Specifications}
\label{sec:robust}

The primary configuration fixes a random seed, a covariance estimation window and a
rebalancing interval. Table~\ref{tab:sweep} reports every combination of three seeds, two
estimation windows and three rebalancing intervals, so that the reported result can be
located within the distribution of specifications rather than presented in isolation.

\begin{table}[!ht]
\centering
\caption{Specification sweep. All eighteen cells reported; the last two columns record
whether the cell beats the equal-weight benchmark on each measure. Sharpe ratios here are
the annualised arithmetic mean over volatility, so that all eighteen cells are computed
identically; on the same basis the equal-weight benchmark attains 0.670 with a maximum
drawdown of $-36.1$\%.}
\label{tab:sweep}
\footnotesize
\begin{tabular}{rrrrrcc}
\hline
\textbf{Seed} & \textbf{Cov. window} & \textbf{Rebalance} & \textbf{Sharpe} & \textbf{Max DD} & \textbf{Sharpe win} & \textbf{DD win} \\
\hline
0 & 126 & 10 & 0.868 & $-20.7$\% & yes & yes \\
0 & 126 & 21 & 0.954 & $-19.6$\% & yes & yes \\
0 & 126 & 42 & 0.924 & $-19.7$\% & yes & yes \\
0 & 252 & 10 & 0.823 & $-20.8$\% & yes & yes \\
0 & 252 & 21 & 0.905 & $-19.0$\% & yes & yes \\
0 & 252 & 42 & 0.899 & $-20.4$\% & yes & yes \\
7 & 126 & 10 & 0.881 & $-20.6$\% & yes & yes \\
7 & 126 & 21 & 0.956 & $-19.6$\% & yes & yes \\
7 & 126 & 42 & 0.912 & $-19.6$\% & yes & yes \\
7 & 252 & 10 & 0.818 & $-21.1$\% & yes & yes \\
7 & 252 & 21 & 0.891 & $-19.0$\% & yes & yes \\
7 & 252 & 42 & 0.888 & $-20.4$\% & yes & yes \\
42 & 126 & 10 & 0.853 & $-20.6$\% & yes & yes \\
42 & 126 & 21 & 0.928 & $-19.6$\% & yes & yes \\
42 & 126 & 42 & 0.893 & $-19.7$\% & yes & yes \\
42 & 252 & 10 & 0.813 & $-21.1$\% & yes & yes \\
42 & 252 & 21 & 0.887 & $-19.0$\% & yes & yes \\
42 & 252 & 42 & 0.853 & $-20.4$\% & yes & yes \\
\hline
\end{tabular}
\end{table}

Every one of the eighteen specifications beats the equal-weighted benchmark on both the
Sharpe ratio and maximum drawdown. The Sharpe ratio ranges from 0.813 to 0.956 and the
maximum drawdown from $-21.1$ to $-19.0$ per cent. The primary configuration appears in
this table at 0.887, which is the median of the eighteen values: it is not the best cell,
and the best cell exceeds it by 0.069. Sensitivity to the random seed is small, which
indicates that the mixture fit is stable; sensitivity to the rebalancing interval is
larger, with ten-session rebalancing consistently weakest because it incurs more turnover
for no additional information.

The same sweep is unfavourable on one comparison. In none of the eighteen cells does the
framework beat the volatility-matched equal-weighted benchmark on maximum drawdown. This
is a consistent finding rather than an artefact of the primary configuration, and it
bounds the drawdown claim precisely: the framework reduces drawdown relative to an
unlevered diversified benchmark, not relative to every portfolio held at the same
volatility.

\subsubsection{External Validation of the Sentiment Composite}
\label{sec:sentval}

Naming the second composite a sentiment measure is a claim about what it measures, and a
claim of that kind is testable without reference to whether it makes money. This
subsection tests it; Section~\ref{sec:ic} asks the separate question of whether it
forecasts anything.

The composite is aggregated to the market level by averaging each of its four inputs
across the ten instruments, standardising each average over the sample and combining the
four with the signs of Table~\ref{tab:features}. The resulting series is compared with
four published measures, none of which is observed anywhere in the framework: the
University of Michigan survey of consumer sentiment, the Cboe volatility index, the
St.~Louis Fed financial stress index and the Chicago Fed national financial conditions
index. The first is a household survey and contains no market data at all, which makes it
the most demanding of the four. A sentiment measure should move with it and against the
remaining three.

\begin{table}[!ht]
\centering
\caption{External validation of the sentiment composite, January 2008 to August 2026.
Newey--West $t$-statistics in parentheses, with twenty-one lags for daily, eight for
weekly and six for monthly data. The change columns correlate first differences and so
are not driven by any common trend in levels.}
\label{tab:sentval}
\footnotesize
\begin{tabular}{p{4.5cm}crr}
\hline
\textbf{Series} & \textbf{Expected} & \textbf{Levels} & \textbf{Changes} \\
\hline
Michigan consumer sentiment, monthly ($n=222$) & $+$ & $+0.262$ ($+1.91$) & $+0.240$ ($+3.54$) \\
\hline
Cboe volatility index, daily ($n=4{,}689$) & $-$ & $-0.541$ ($-3.62$) & $-0.112$ ($-3.02$) \\
\hline
St.~Louis Fed financial stress index, weekly ($n=972$) & $-$ & $-0.606$ ($-2.71$) & $-0.254$ ($-3.68$) \\
\hline
Chicago Fed financial conditions index, weekly ($n=972$) & $-$ & $-0.513$ ($-2.46$) & $-0.123$ ($-2.18$) \\
\hline
\end{tabular}
\end{table}

Every correlation in Table~\ref{tab:sentval} carries the expected sign, and every one of
the four is significant in first differences at the five per cent level. The survey
measure is the informative case. A correlation of $+0.240$ between monthly changes in a
composite built from exchange turnover and monthly changes in what households tell a
telephone interviewer is not a mechanical artefact of shared construction, because the
two series share no inputs; it is evidence that the composite is picking up the quantity
its name asserts. The financial stress and financial conditions indices, which do contain
market data, correlate more strongly in levels at $-0.606$ and $-0.513$.

The relationship is stable in sign though not in magnitude. Splitting the sample into
four disjoint blocks---2008--2013, 2014--2019, 2020--2023 and 2024--2026---the sign is
correct in all sixteen series-by-block combinations. Magnitudes are largest in the two
blocks containing crises, reaching $-0.823$ against the stress index over 2008--2013, and
smallest over the calm 2014--2019 block, where the correlation with that index falls to
$-0.196$. This is the expected pattern rather than a defect: a measure of stress has
little to discriminate when there is little stress. Replacing the full-sample
standardisation with an expanding-window standardisation that uses only past data leaves
every sign unchanged, at $+0.182$, $-0.341$, $-0.406$ and $-0.319$ respectively.

The second test bears directly on the framework's central mechanism. If the identified
market states are states of sentiment and volatility, the composite should differ across
them. Table~\ref{tab:sentregime} evaluates the composite at the 223 rebalance dates for
which a regime label exists.

\begin{table}[!ht]
\centering
\caption{Sentiment composite by identified market state, measured at the rebalance dates.
The states are labelled by fitted volatility, not by sentiment, and the sentiment
composite plays no part in fitting them.}
\label{tab:sentregime}
\footnotesize
\begin{tabular}{p{4.5cm}rrr}
\hline
\textbf{Identified state} & \textbf{$n$} & \textbf{Mean} & \textbf{SD} \\
\hline
Calm & 160 & $+0.073$ & 0.407 \\
\hline
Transitional & 56 & $-0.140$ & 0.382 \\
\hline
Stress & 7 & $-0.831$ & 0.735 \\
\hline
\end{tabular}
\end{table}

The ordering is monotone and steep, and the differences are unlikely to be accidental:
one-way analysis of variance gives $F = 19.34$ with $p = 1.8 \times 10^{-8}$, and the
rank-based Kruskal--Wallis test, which does not assume normality and is the safer of the
two given seven observations in the stress state, gives $H = 29.54$ with
$p = 3.9 \times 10^{-7}$. The mixture model is fitted on realised volatility and
short-horizon momentum; the sentiment composite is not among its inputs. That the states
separate this sharply on a variable withheld from their construction is what justifies
describing them as sentiment and volatility states rather than volatility states alone.

Two limits should be read alongside these results. The stress state contains seven
rebalance dates, so its mean is estimated imprecisely even though the test statistic is
large. And nothing established here concerns forecasting. A composite can measure
sentiment accurately and still fail to predict returns, which is precisely what
Section~\ref{sec:ic} documents; validity and profitability are separate properties, and
conflating them is a common error in this literature.

\subsubsection{Component Ablation}
\label{sec:ablation}

The framework has two active layers, signal fusion and the regime risk budget.
Table~\ref{tab:ablation} removes each in turn. The comparison requires care, because
scaling a portfolio toward cash changes its return and its volatility together; comparing
a de-risked strategy against a fully invested variant therefore confounds the timing of
risk with its level. The first ablation holds average invested exposure fixed at the
framework's own 92.1 per cent and applies it constantly, so that only the timing differs.

\begin{table}[!ht]
\centering
\caption{Component ablation. The matched-exposure control holds the same average
invested exposure as the full framework but applies it without regard to the regime.}
\label{tab:ablation}
\footnotesize
\begin{tabular}{p{4.8cm}rrrrr}
\hline
\textbf{Configuration} & \textbf{CAGR} & \textbf{Sharpe} & \textbf{Max DD} & \textbf{Exposure} & \textbf{$p$} \\
\hline
Full framework & 5.65\% & 0.878 & $-18.98$\% & 92.1\% & --- \\
No regime timing, same average exposure & 5.43\% & 0.795 & $-22.26$\% & 92.1\% & 0.695 \\
No regime timing, full exposure & 6.91\% & 0.923 & $-22.53$\% & 100.0\% & 0.016 \\
No signal layer & 5.75\% & 0.908 & $-18.87$\% & 92.1\% & 0.650 \\
Equal-weight benchmark & 8.45\% & 0.628 & $-36.08$\% & 100.0\% & --- \\
\hline
\end{tabular}
\end{table}

The first comparison is the one the framework is built on, and it succeeds. Holding
average exposure fixed, allocating that exposure according to the regime rather than
uniformly raises the Sharpe ratio from 0.795 to 0.878 and improves maximum drawdown from
$-22.26$ to $-18.98$ per cent. The difference in mean return is not significant
($p = 0.695$), which is the expected signature of a pure risk-timing effect: the same
average risk is carried, but it is carried at better moments. This is the only comparison
in the study in which the framework improves both measures at once, and it is the claim
the paper rests on.

The second comparison bounds the cost of that timing. Against a fully invested variant,
the framework gives up approximately 1.3 percentage points of annual return, and the
difference is statistically significant ($p = 0.016$). De-risking is not free, and its
Sharpe ratio under this comparison is lower, 0.878 against 0.923. What the framework buys
with that return is 3.55 percentage points of maximum drawdown. Whether that exchange is
worthwhile is a preference, not a fact, and the paper reports both sides of it.

The third comparison is a null, and it concerns the layer the framework was named for.
Removing the signal fusion layer entirely---setting expected returns to zero so that the
optimiser becomes a constrained minimum-variance problem---changes the Sharpe ratio from
0.878 to 0.908 and the maximum drawdown from $-18.98$ to $-18.87$ per cent. Both
movements favour removal, and neither is statistically significant ($p = 0.650$). The
signal fusion layer contributes nothing detectable to portfolio outcomes in this
universe, and its point estimate is mildly adverse. Section~\ref{sec:ic} establishes why.

\textbf{Form of the Risk Budget:}

Table~\ref{tab:budget} varies the risk budget itself, including the absolute-target
formulation used in the framework as originally specified.

\begin{table}[!ht]
\centering
\caption{Risk budget specifications. Relative budgets are stated as the fraction of the
sleeve's own volatility retained in the calm, transitional and stress states.}
\label{tab:budget}
\footnotesize
\begin{tabular}{p{4.6cm}rrrr}
\hline
\textbf{Budget} & \textbf{CAGR} & \textbf{Sharpe} & \textbf{Max DD} & \textbf{Exposure} \\
\hline
Relative 1.00 / 0.75 / 0.50 (primary) & 5.65\% & 0.878 & $-18.98$\% & 92.1\% \\
Relative 1.00 / 0.80 / 0.60 & 5.77\% & 0.871 & $-19.95$\% & 93.7\% \\
Relative 1.00 / 0.60 / 0.30 & 5.31\% & 0.881 & $-16.62$\% & 87.6\% \\
Relative 1.00 / 0.50 / 0.25 & 5.21\% & 0.891 & $-15.71$\% & 84.9\% \\
Absolute targets, 20 / 15 / 10\% annualised & 6.93\% & 0.923 & $-22.86$\% & 100.0\% \\
\hline
\end{tabular}
\end{table}

Two findings follow. First, the Sharpe ratio varies little across the four relative
budgets, from 0.871 to 0.891, while maximum drawdown improves monotonically as the budget
tightens, from $-19.95$ to $-15.71$ per cent. The result is therefore not an artefact of
the particular multipliers chosen; the mildest of the four is reported as primary
precisely because it is the least aggressive, and the tighter budgets would have produced
a marginally better headline on both measures.

Second, the absolute-target formulation is inert. Its mean invested exposure is 99.99 per
cent, and the constraint binds in one rebalance out of 224. Targets of 10 to 20 per cent
annualised volatility are never reached by a sleeve whose own realised volatility is 6.4
per cent, so under that formulation the regime layer is present in the specification but
has no effect on the portfolio. Its apparently superior Sharpe ratio of 0.923 is simply
the fully invested portfolio of Table~\ref{tab:ablation}, reached by a different route.
This is the practical reason the relative formulation is adopted: an absolute volatility
target is not scale-free, and a target calibrated for one universe silently disables the
mechanism in another.

\subsubsection{Statistical Significance}

Table~\ref{tab:inference} reports inference on the daily excess return series and on the
Sharpe difference.

\begin{table}[!ht]
\centering
\caption{Inference against the equal-weighted benchmark. Bootstrap intervals are from a
stationary block bootstrap with 21-session blocks and 2{,}000 replications.}
\label{tab:inference}
\footnotesize
\begin{tabular}{p{6.4cm}r}
\hline
Mean excess return over benchmark, annualised & $-3.31$\% \\
Paired $t$-statistic & $-1.369$ \\
$p$-value & 0.171 \\
Newey--West $t$-statistic, ten lags & $-1.699$ \\
\hline
Strategy Sharpe ratio, 95\% bootstrap interval & $[0.394, \; 1.368]$ \\
Benchmark Sharpe ratio, 95\% bootstrap interval & $[0.180, \; 1.165]$ \\
Sharpe difference, median & $+0.242$ \\
Sharpe difference, 95\% bootstrap interval & $[-0.102, \; 0.574]$ \\
Bootstrap probability that the difference is not positive & 0.088 \\
\hline
\end{tabular}
\end{table}

The results must be stated carefully. The framework's mean return is below the
benchmark's, and that shortfall is not statistically significant ($p = 0.171$). The
Sharpe difference of $+0.250$ has a bootstrap median of $+0.242$ and a 95 per cent
interval of $[-0.102, 0.574]$ that includes zero; the bootstrap probability that the
difference is not positive is 0.088. On an eighteen-year sample, therefore, the Sharpe
improvement is not significant at conventional levels, and the paper does not claim that
it is. Extending the sample by the 665 sessions of the holdout narrowed that
interval---its lower bound moves from $-0.194$ on the pre-2024 data to $-0.102$---without
crossing zero, which is the expected behaviour of a real but modest effect measured on a
sample that remains too short.

What is robust is not the point estimate but its consistency. The Sharpe advantage
appears in 18 of 18 specifications, the drawdown advantage in 18 of 18 and in 18 of 19
calendar years, and the outperformance in all four stress episodes. Eighteen years of
daily data contain few independent observations of the events that matter for a drawdown
claim---four, on the evidence of Table~\ref{tab:episodes}---and no amount of daily
sampling converts that into statistical power \citep{harvey15}. The appropriate
conclusion is that the drawdown reduction is consistent and mechanically explicable, and
that the Sharpe improvement, while positive under every specification examined, is not
established at conventional significance.

\subsubsection{Secondary Universe: Mega-Capitalisation Equities}
\label{sec:megacap}

The framework was also run without modification on the ten mega-capitalisation US
equities of the original study, over the same 2018--2026 download window, giving 1{,}920
evaluated sessions from January 2019. Table~\ref{tab:megacap} reports the outcome, which
is unfavourable on risk-adjusted return.

\begin{table}[!ht]
\centering
\caption{Secondary universe: ten mega-capitalisation US equities, 2019--2026}
\label{tab:megacap}
\footnotesize
\begin{tabular}{p{4.8cm}rrrr}
\hline
& \textbf{CAGR} & \textbf{Volatility} & \textbf{Sharpe} & \textbf{Max DD} \\
\hline
Regime-conditional strategy & 17.02\% & 14.87\% & 1.145 & $-20.45$\% \\
Equal-weight benchmark & 33.52\% & 24.45\% & 1.371 & $-33.63$\% \\
Equal-weight at strategy volatility & 20.08\% & 14.87\% & 1.350 & $-21.41$\% \\
No regime timing, same average exposure & 18.82\% & 16.82\% & 1.119 & $-30.36$\% \\
\hline
\end{tabular}
\end{table}

The framework loses on risk-adjusted return, 1.145 against 1.371, and gives up 15.09
percentage points of annual return, a shortfall that is statistically significant
($t = -2.791$, $p = 0.005$). Its defensive behaviour nevertheless survives intact:
maximum drawdown improves from $-33.63$ to $-20.45$ per cent, it wins on drawdown in 8 of
8 calendar years while winning on Sharpe ratio in only 3, and it outperforms in both
stress episodes this window contains, by 12.50 percentage points during the 2020 crash
and 12.60 percentage points during the 2022 bear market.

The regime layer itself continues to behave as the primary results describe. Against the
matched-exposure control it improves the Sharpe ratio slightly, 1.145 against 1.119, and
improves maximum drawdown substantially, from $-30.36$ to $-20.45$ per cent. The
risk-timing effect identified in Section~\ref{sec:ablation} therefore reproduces in a
universe on which the framework loses overall, which is the relevant test of whether that
effect is specific to the primary sample. What fails here is not the regime layer but the
decision to hold the resulting portfolio in place of the benchmark.

The reason is a property of the universe rather than of the framework. Ten
mega-capitalisation US equities are close to a single factor bet: correlations are high,
there is no low-volatility asset to rotate into, and the risky sleeve therefore carries
roughly 15 per cent annualised volatility against 6 per cent in the multi-asset case.
De-risking in such a universe can only move capital to cash, so the framework
systematically forgoes return during an exceptional advance and recovers only part of it
in drawdowns. This bounds the applicability of the framework directly: it requires a
universe with genuine dispersion in asset volatility, and it should not be expected to
add value in a concentrated, single-sector one. The originally submitted version of this
study used exactly such a universe, which is why its results are superseded here.

\subsection{Theoretical and Practical Implications}

\subsubsection{Methodological Contributions}

\begin{itemize}
\item \textbf{Signal-allocation bridge.} The research demonstrates a systematic framework
for translating predictive signals into executable portfolio decisions, addressing a
critical gap in the literature. What it additionally shows is that the bridge can be built
and the signal crossing it can still be uninformative, which is why each layer is tested
separately below.
\item \textbf{Regime adaptation.} The regime-dependent risk budget improves both
risk-adjusted return and drawdown at matched average exposure. The regime-dependent
\emph{weighting} mechanism, by contrast, is not supported: the conditional efficacy it
presumes is not present in the data.
\item \textbf{Practical implementation.} The incorporation of transaction costs and
realistic constraints ensures the approach is implementable in live trading environments,
and the cost assumption is shown to act as a regulariser rather than merely as a
deduction.
\end{itemize}

The framework's contribution is therefore narrower than the one claimed for it at the
outset, and better supported. Three findings survive the evaluation above. First,
splitting the problem into relative composition and total risk, rather than letting one
constraint set determine both, makes each part independently testable, and is what allows
the ablation in Table~\ref{tab:ablation} to be constructed at all. Second, regime
conditioning earns its place through timing rather than level: at identical average
exposure it improves the Sharpe ratio from 0.795 to 0.878 and maximum drawdown from
$-22.26$ to $-18.98$ per cent with no significant difference in mean return
($p = 0.695$), at a measured cost of approximately 1.3 percentage points of annual return
against a fully invested variant, significantly so, for 3.6 percentage points of maximum
drawdown. Third, the form of the risk budget matters more than its calibration---the four
relative budgets in Table~\ref{tab:budget} produce Sharpe ratios within 0.020 of one
another, while the absolute-target formulation, identical in intent, is inert, binding in
one rebalance out of 224. An absolute volatility target is not scale-free, so a
specification moved between universes without re-calibration can silently disable the
mechanism it was written to implement.

\textbf{What the Results Do Not Establish:}

Two claims the framework was designed to support are not supported. The signal fusion
layer does not contribute: removing it improves the Sharpe ratio from 0.878 to 0.908 and
leaves maximum drawdown unchanged, neither movement significant, and
Table~\ref{tab:ic} explains why---both composites are significantly negatively related to
next-session cross-sectional returns over the full sample. The regime-conditional
weighting schedule has no measured basis either. It presumes that sentiment information
becomes relatively more informative under stress; the point estimates run the other way
and the calm-to-stress contrast is not significant ($t = -1.46$, $p = 0.145$). Split at
the holdout boundary, the fused coefficient is $-0.039$ before 2024 and $+0.009$ after,
so even the negative relationship is era-specific and could not have been exploited by
inverting the composites. Nor is the improvement in risk-adjusted return statistically
significant: the Sharpe difference of $+0.250$ carries a bootstrap interval of
$[-0.102, 0.574]$, and adding the 665 sessions of the holdout narrowed that interval
without moving it clear of zero. Consistency across 18 of 18 specifications, 18 of 19
calendar years and 4 of 4 episodes is evidence of robustness rather than of significance.
The two should not be conflated, and this paper does not conflate them.

\textbf{Implications for Practice:}

A regime-conditional risk budget is a reasonable component of a multi-asset allocation
process whose objective includes drawdown control and which can sacrifice some return for
it. It should be implemented in relative rather than absolute form, evaluated against a
volatility-matched benchmark rather than an unlevered one, and not expected to help in a
universe without genuine dispersion in asset volatility. The results also caution against
a common design pattern: the framework as originally specified contained two layers whose
contribution had never been isolated, one inert by construction and the other pointed in
the wrong direction, and neither defect was visible in aggregate performance
\citep{prado16}.

\subsubsection{Limitations and Future Research}
\label{sec:future}

Several limitations bound what can be concluded.

\begin{itemize}
\item The analysis relies on simplified proxy variables for market sentiment, foregoing
the capabilities of contemporary large language model processing of textual data.
\item The weighting parameters applied within the regime-dependent fusion mechanism are
statically defined, lacking an adaptive framework for their continuous optimisation.
\item The empirical validation is confined to a constrained asset universe and a specific
historical timeframe, potentially limiting the generalisability of the findings.
\end{itemize}

Each bears on the results in a specific way. The universe is small: ten exchange-traded
funds demonstrate the mechanism but do not establish that it generalises to larger or
differently composed universes, and the 25 per cent position cap binds on average, so the
optimiser's freedom is partly set by the constraint rather than by the risk model. The
evaluation is also a single historical path---eighteen years is long by the standards of
this literature and short for inference about tail events, and the four stress episodes
are not independent draws from a stationary process.

The signals are market-derived proxies rather than textual sentiment: no news, filings or
social media text is used at any point, as stated in Section~\ref{sec:sentiment}. The
null reported in Section~\ref{sec:ic} is therefore a null about these proxies, not
evidence that textual sentiment would fail in the same role; establishing that would
require a time-stamped news corpus covering the traded instruments, which for broad
exchange-traded funds does not exist at the instrument level. The regime model is
deliberately simple---a three-component Gaussian mixture on two features is interpretable
and stable, but imposes conditional normality within states, treats observations as
conditionally independent given the state, and, as Table~\ref{tab:regimefit} shows, is
not the specification either information criterion would select. Finally, the risk-free
rate is taken as zero and cash assumed to earn nothing; over a sample spanning both the
near-zero rates of 2009--2015 and the higher rates of 2023 this understates the return on
the framework's cash holding and so understates its performance against the fully
invested benchmark. That bias runs against the framework, which is the conservative
choice, but it is a simplification.

\textbf{Research Extensions:}

Four extensions follow directly from the limitations above.

\begin{itemize}
\item \textbf{Advanced sentiment integration.} A logical extension involves the
incorporation of language-model-driven sentiment analysis to derive more sophisticated
indicators of market affect, applied to a universe where an instrument-level news corpus
exists. This would convert the proxy null of Section~\ref{sec:ic} into a genuine test of
the fusion hypothesis.
\item \textbf{Dynamic optimisation frameworks.} The implementation of reinforcement
learning algorithms could facilitate the dynamic calibration of regime-specific
parameters, moving beyond the current static configuration. This is subject to the
overfitting concerns documented by \citet{harvey15} and \citet{prado16}.
\item \textbf{Expanded portfolio scope.} The framework's robustness could be further
tested through application to global, multi-asset portfolios encompassing equities, fixed
income, commodities and currencies. Such a test would establish whether the
volatility-dispersion condition identified in Section~\ref{sec:megacap} is the binding
requirement it appears to be.
\item \textbf{Alternative data synthesis.} Future work could enhance the feature set by
integrating a broader spectrum of alternative data, such as supply chain information or
electronic market microstructure data, which would widen the test of the fusion hypothesis
beyond the price and volume aggregates used here.
\item \textbf{Alternative regime formulations.} Replacing the mixture model with a
formulation that represents state persistence explicitly, such as a hidden Markov or jump
model \citep{li25,luo26}, would remove the conditional-independence assumption noted
above.
\end{itemize}
